\section{The Field You Cannot Ship}
\label{sec:ch01-the-field-you-cannot-ship}

\index{coordination problem|(}%
Booking.com built its GraphQL layer the way many teams build their first one:
a single service that called the services behind it and mapped their responses
into one schema. Christian Ernst, writing on Booking.com's engineering blog,
wrote down what that had cost by the time it grew up. \enquote{As our single
GraphQL service grew, it became more difficult to maintain not only for the
GraphQL team who had dozens of merge requests a day to review for schema
changes and service mappings but also manage frequent service deployments and
help resolve merge conflicts with multiple teams}~\autocite{ernst2022}.

Read that again for who is doing the work. Every schema change in the company
arrives as a merge request in front of one team, which reviews it, deploys it,
and untangles it when two other teams collide in the same file. The teams that
own the data are not the teams shipping the change, and Ernst names what
follows: \enquote{This also made it difficult for service owners to have strong
ownership of their data and schema as they were heavily tied to the single
service}~\autocite{ernst2022}.

Netflix arrived at the same place from a different direction. Its Studio API
was a curated graph API over the data behind Netflix's studio applications, and
Tejas Shikhare's account of it is unusually honest, because the thing worked
first. \enquote{The One Graph exposed by Studio API was a runaway success;
product teams loved the reusability and easy, consistent data
access}~\autocite{shikhare2020}. Then it got big, and three bottlenecks
appeared. Two of them are this chapter's subject: \enquote{the Studio API team
was disconnected from the domain expertise and the product needs, which
negatively impacted the schema's health}, and \enquote{it was hard for one
small team to handle the increasing operational and support burden for the
expanding graph}~\autocite{shikhare2020}.

Neither company describes a technical failure. In both, one team sits between
every other team and the schema, and that team is the bottleneck when it is
busy and the wrong team to ask when it is not, because it does not own the
domains it is being asked to model.

Figure~\ref{fig:ch01-one-field} traces one change through both arrangements:
you need a single field added to a type another team owns. Nobody in lane A is
careless. Your change is scheduled against work you do not control and cannot
see.
\index{coordination problem|)}%

\begin{figure}
  \centering
  % Adding one field, in a shared graph and in a federated one.
%
% This is the book's first figure and therefore fixes the idiom SPEC decision
% 33 describes: square corners, no fills, one stroke weight, sans-serif labels,
% black on white, and order shown as a timeline rather than as a box-and-arrow
% architecture diagram. Later figures are matched against this file, not
% against the decision's wording.
%
% Bare tikzpicture (house style): the figure environment, caption and label
% live at the call site in the section file.
\begin{tikzpicture}[
  x=1mm, y=1mm,
  every node/.append style={font=\sffamily\scriptsize},
  axis/.style={draw, line width=0.4pt, -{Straight Barb[length=1.4mm, width=1.6mm]}},
  tick/.style={draw, line width=0.4pt},
  event/.style={anchor=south, align=center, text width=33mm, inner sep=1pt},
  span/.style={draw, line width=0.4pt},
  spanlabel/.style={anchor=north, align=center, inner sep=2pt},
  lane/.style={anchor=west, font=\sffamily\scriptsize\itshape},
]

% ------------------------------------------------------------------- lane A
\node[lane] at (0,13) {A. one schema, owned by one team};

\draw[axis] (0,0) -- (150,0);
\foreach \x in {15,55,95,135}{\draw[tick] (\x,-1.5) -- (\x,1.5);}

\node[event] at (15,3)  {you ask for\\the field};
\node[event] at (55,3)  {it joins\\their queue};
\node[event] at (95,3)  {they edit the schema\\and the resolver};
\node[event] at (135,3) {one deployment,\\everyone's};

\draw[span] (48,-4) -- (48,-7) -- (102,-7) -- (102,-4);
\node[spanlabel] at (75,-8) {another team's calendar, and you are not on it};

% ------------------------------------------------------------------- lane B
\begin{scope}[shift={(0,-42)}]
  \node[lane] at (0,13) {B. a federated graph};

  \draw[axis] (0,0) -- (150,0);
  \foreach \x in {20,75,130}{\draw[tick] (\x,-1.5) -- (\x,1.5);}

  \node[event] at (20,3)  {you edit your\\own subgraph};
  \node[event] at (75,3)  {it composes,\\or it does not};
  \node[event] at (130,3) {your deployment,\\yours alone};

  \draw[span] (64,-4) -- (64,-7) -- (86,-7) -- (86,-4);
  \node[spanlabel] at (75,-8) {a machine, and it answers every time you ask};
\end{scope}

\end{tikzpicture}

  \caption{One new field on a type somebody else owns, in each arrangement.
  Lane A waits on another team's priorities, which you cannot query. Lane B
  waits on a program, which answers the same way every time you run it. That
  difference, rather than the number of steps, is what federation buys.}
  \label{fig:ch01-one-field}
\end{figure}

This is worth being precise about, because it is the test that decides whether
the rest of the book is for you. The pain is not that the schema is large. It
is that the number of teams who need to change it exceeds the number of teams
who are allowed to.
